L'usage de l'ordinateur se généralise dans tous les domaines de l'activité
économique. La plupart des utilisateurs emploient des logiciels 
d'ap\-pli\-ca\-tion
pour accomplir une ou plusieurs tâches spécifiques~: édition de
texte, comp\-ta\-bi\-li\-té, jeux, etc. Par désir de tirer tout le parti de la 
puissance de la machine, par plaisir ludique ou par simple
envie de comprendre le principe de fonctionnement des ordinateurs, de
nombreux informaticiens en herbe ont envie d'apprendre, parfois seuls, la
programmation. La tâche se révèle souvent longue, difficile
voire rébarbative. Nous voudrions en faire une activité créatrice,
formatrice pour l'esprit et passionnante.

\medskip

De nombreuses introductions à la programmation reposent sur la connaissance
préalable de l'architecture des ordinateurs ou utilisent des notions 
mathématiques qui ne sont pas à la portée des enfants ou des adultes
n'ayant pas une formation scientifique \cite{LT}\@%
\footnote{
Les scientifiques ont l'habitude de `citer leurs sources' en
donnant les références aux articles de revues scientifiques et aux 
livres des auteurs dont ils utilisent les idées. La méthode moderne pour
citer les sources est d'utiliser un numéro placé entre crochets 
(comme par exemple \cite{LT}) qui renvoie à une liste de 
références, placée à la fin du livre, que l'on appelle 
{\em bibliographie} (voir page 
\pageref{ch:biblio}).}. {\footnotesize (Les petits 
chiffres surélevés renvoient à une note de
 bas de page expliquant le mot qui précède. Par exemple, la note qui 
 se trouve en bas de cette page explique l'usage des chiffres entre crochets 
 comme \cite{LT}.)} Nous proposons de procéder autrement.

\medskip

Nous uti\-li\-sons le dessin, qui est à la portée de tous. Lors de la première publication du livre
en 1991, cette approche était
relativement nouvelle dans l'apprentissage de l'informatique, les progrès 
sai\-sis\-sants de la technologie qui ont permis de mettre à la portée du 
grand public des micro-ordinateurs dotés d'écrans graphiques très 
performants ne datant que des années 1980. Il est plaisant de noter que
les origines de notre activité graphique remontent à la préhistoire (plus précisément à
la période de l'art pariétal%
%
\footnote{Peinture sur les parois d'une grotte.
} entre 22000 et 13000 ans av. J.-C.). L'informatique graphique qui est née pour le grand public
dans les années 1980 va donc nous permettre à 20000 ans d'intervalle de dessiner avec 
le même plaisir que nos lointains ancêtres. 

\medskip

Ce livre s'adresse à tous. Les enfants des classes maternelles sont capables
de faire des dessins stylisés à la main et avec un peu d'imagination, de les 
reproduire sur l'écran
de l'ordinateur en utilisant un programme de dessin interactif (chapitres 1 et 2).
Les enfants des classes primaires peuvent écrire des programmes pour réaliser 
des dessins  figuratifs (chapitres 3 à 5). Très vite ils découvriront la facilité avec 
la\-quel\-le ils peuvent programmer l'ordinateur pour faire des dessins géométriques simples basés sur des
motifs répétitifs qu'ils n'auraient pas la patience ni la minutie d'exécuter
à la main (chapitres 6 et 7). Les collégiens mettront en œuvre leurs
connaissances rudimentaires d'algèbre et de géométrie analytique pour écrire des 
programmes de dessin
utilisant des calculs élémentaires et des repères cartésiens (chapitres 8 à 17). 
Ils trouveront une base 
intuitive pour mettre en œuvre des
concepts qui peuvent paraître abstraits. Les
lycéens sont à même d'assimiler des notions plus informatiques comme le 
codage, voire des notions moins élémentaires comme la 
récursivité, qui deviennent évidentes par le dessin (chapitres 17 à 22). Les étudiants
trouveront du plaisir à comprendre simplement la programmation qui paraît parfois obscure quand
elle est présentée trop rapidement en premier cycle universitaire. 

\medskip

Les langages de programmation n'offrant pas les possibilités graphiques 
souhaitables, nous avons ajouté une bi\-blio\-thè\-que à \pascal\@%
%
\footnote{Nous utiliserons les 
compilateurs Turbo Pascal 5.5 ou 6.0 de Borland sur les or\-di\-na\-teurs compatibles 
IBM PC et le compilateur Turbo Pascal 1.1 de Borland sur les Macintosh.}
%
qui nous permet
d'obtenir de façon simple des effets simples (comme les arcs de cercles) ou
spectaculaires (comme la peinture) sans être liés à la technologie des
écrans graphiques actuels et sans faire appel à des notions avancées de
programmation comme les structures de données. Cette bibliothèque ainsi
que tous les exemples et programmes solutions des exercices proposés dans le livre
sont fournis sur disquette%
%
\footnote{\ Un bon de commande de la disquette d'accompagnement se
trouve à la fin du livre. Pour apprendre à l'utiliser, l'insérer dans un 
lecteur puis, sur compatible IBM PC taper {\tt LisezMoi} ou sur Macintosh
cliquer deux fois rapidement avec la souris sur l'icône {\tt LisezMoi}.}%
. Elle offre de nombreuses possibilités qui ne sont pas explorées dans le livre,
comme par exemple pour les spécialistes, de créer des programmes interactifs d'apprentissage
comme ceux qui sont proposés pour l'addition, les coordonnées cartésiennes et la
translation.

\medskip

Ce livre devrait être lu par étapes successives, en prenant le temps et le plaisir 
de faire de nombreuses expériences personnelles sur ordinateur, en résolvant des exercices,
en modifiant les programmes proposés en solution et en
recherchant de nouveaux thèmes de dessin.


\vspace{1.0cm}

%%%%%%%%%%%%%%%%%%%%%%%%%%%%%%%%%%%%%%%%%%%%%%%%%%%%%%%%%%%%%%%%%%%%%%%%%%%
\noindent \framebox[\textwidth]{
\begin{minipage}{\myparboxwidth}
\setlength{\parindent}{2em}
\small \sf
\vspace{\abovedisplayskip}
\section*{Addenda}

Nous donnons en addenda à la fin de chaque chapitre des notes additionnelles
destinées aux lecteurs déjà familiers de l'informatique.

Nous introduisons dans ce livre les activités de base de la programmation
comme la saisie du texte, l'interprétation, la compilation, l'exécution
et la mise au point de programmes. Les notions de programmation sont présentées
en \pascal~: structuration des programmes en procédures, itérations
bornées, expressions entières, déclarations de constantes, 
paramètres (passés par valeur), tests, expressions booléennes,
itérations non bornées, codages de l'information, expressions rationnelles 
(réelles), récursivité, variables de types simples et affectations. 
Notre approche basée sur le dessin permet de donner une intuition 
graphique à ces notions et évite les problèmes plus complexes de
spécification et de correction des programmes (qui sont des évidences 
visuelles). Les structures de données (tableaux,
enregistrements et pointeurs en \pascal) sont volontairement laissées pour
des apprentissages ultérieurs plus élaborés.

L'apprentissage de la programmation à un niveau très élémentaire comme
celui qui est proposé dans ce livre permet d'acquérir le sens du passage du concret 
à l'abstrait, l'intuition des structures formelles, le raisonnement inductif et la capacité à 
formaliser qui sont à la 
base de la programmation
mais également le point de départ de toute démarche scientifique.

\vspace{\abovedisplayskip}
\end{minipage}
}
%%%%%%%%%%%%%%%%%%%%%%%%%%%%%%%%%%%%%%%%%%%%%%%%%%%%%%%%%%%%%%%%%%%%%%%%%%%%
